\documentclass[english]{syllabus}

\coursecode{ICS0036}
\coursetitle{Cryptography}
\semester{autumn}{2026}

\begin{document}

\makesyllabustitle

\field{objectives}{This course aims to introduce fundamental cryptographic primitives, techniques and protocols with a focus on their applications, and to build students' ability to design and analyse secure digital systems.}

\field{outcomes}{
  After completing this course, the student:
  \begin{itemize}
    \item describes the basic properties, functions and limitations of common cryptographic primitives;
    \item explains how cryptographic protocols (e.g. TLS) combine different primitives to achieve multiple security goals, and identifies and explains basic cryptographic issues in a protocol design;
    \item knows which cryptographic algorithms and protocols to use in practice, how to securely instantiate them, and the risks associated with selecting nonstandard options;
    \item understands and justifies the principles of trust networks and PKI;
    \item writes programs/scripts that perform common cryptographic operations using standard libraries and tools;
    \item consults authoritative sources on cryptographic best practices.
  \end{itemize}
}

\field{description}{Pseudorandomness; symmetric and asymmetric cryptosystems; hash functions, MACs and KDFs; digital signatures, certificates, PKI and eIDAS; TLS; side-channel attacks and basics of secure programming; commitment and identification schemes; zero-knowledge proofs and Sigma protocols; secret sharing and secure multi-party computation; online voting; post-quantum cryptography.}

\field{language}{English}

\field{eap}{6 ECTS}

\field{students}{Elective course for the Cybersecurity Engineering programme.}

\stdfield{specialneeds}

\stdfield{registration}

\field{prerequisites}{Familiarity with discrete mathematics and basic programming is expected.
Familiarity with Python is helpful, but not required.}

\field{resources}{Access to a computer and an Internet connection.
During the course you will need to be able to write and run Python $\ge 3.11$.
We will be installing free software including OpenSSL and GMP.
Access to a \emph{*nix} operating system is thus recommended (WSL should be sufficient on Windows).}

\stdfield{lecturers}

\field{contact}{Contact me either by email or on Discord (preferred).
I will generally reply as soon as I can; if you receive no response within 48 hours, please remind me again.}

\field{timetable}{
  Lectures are on Tuesdays, 12:00--13:30.\\
  Practicals are on Wednesdays, 16:00--17:30.

  If I have to cancel a class (e.g. due to sickness or a conference), I will provide a recording and/or a written worksheet to make up for the lost session.
}

\field{process}{
  Participation in both lectures and practicals is recommended, but not mandatory.
  Participating in at least 8 lectures and 8 practicals relaxes your final exam pass requirement (see the \enquote{Assessment criteria for the exam} section).

  Lectures and practicals will be recorded with Echo360.
  However, some recordings may be incomplete if I have to edit them, e.g. due to the Chatham House Rule.

  Lectures will give theoretical foundations while practicals will supplement the theory with further examples and showcase some tools and how to use them.
  Each practical will have some tasks that you must complete.
  You will have two weeks to complete and submit those tasks.

  The material compounds quickly: each topic ties into the listed learning outcomes and cannot be treated in isolation.
  Your workload will explode if you do not learn the material consistently and on time -- beware!

  The recommended way to follow this course is to attend the lectures and practicals, and note down only the \enquote{eureka} moments and biggest pain points.
  Then, preferably the same evening, but no later than the next lecture, you should rewatch the lecture at your own pace, and pause and rewind where necessary.
  You are unlikely to grasp a lecture's material solely by attending it.

  Homework assignments cover groups of topics.
  Solutions must be committed to TalTech's GitLab, but grades and completion are tracked in Moodle.

  There is a midterm (online, in Moodle) that serves as a mandatory self-check.
  Its grade does not affect your final grade, provided that you get at least one point on it.
  Optional, ungraded self-checks will be available for some, if not all, of the weekly topics.

  The course ends with a mandatory exam.
  The main form of the exam is in-person and written (on paper), but there will be limited slots for an oral exam alternative for those who prefer that form.
  The slots are allocated on a \enquote{first-come, first-served} basis, with priority for students who must take the oral exam due to special needs.
}

\field{esupport}{
  \url{https://moodle.taltech.ee/course/view.php?id=37761}

  Moodle also has a link to the course's Discord server.
  Joining the server is not mandatory, but I strongly recommend it.
}

\field{literature}{
  Listed and linked in Moodle.
  All course materials are also available on GitHub at \href{https://github.com/takakv/ics0036-cryptography}{takakv/ics0036-cryptography}.
}

\field{assessment}{
  The course is graded on a 100-point scale.
  \begin{itemize}
    \item 30p are earned from 5 homework assignments, worth 6p each.

    \item 20p are earned from practical tasks.
    The grade will be the percentage of correct submissions.
    For example, assuming that there are 16 graded tasks, correctly submitting 12 of them will yield 15p.
    
    \item 50p are earned from the final exam.
    The exam grade will be scaled to 50p.
  \end{itemize}

  15/30 homework points are necessary, but not sufficient, to pass the course.
  You will have at least 14 days from each assignment's publication date to submit it.
  You can submit assignments up to 7 days late for half of the original points.
  Past the 7 days, the assignment receives 0 points.
  If circumstances prevent you from meeting a deadline, write to me with your reason no later than 4 days before it, and I will consider an extension.

  Similarly, 10/20 practical points are necessary, but not sufficient, to pass the course.
  You will have at least 14 days to submit the graded tasks.
  The tasks are graded automatically, and you have 3 attempts in total for each.
  If your third submission still does not pass, the task will be counted as failed.
  If you are facing issues, ask questions before your third submission.
  A failed task is final unless you can show that the grading was mistaken.
}

\field{exam}{
  The written exam contains a true/false (T/F) section with 20 questions and six sections with short-form open-ended questions:
  \begin{enumerate}
    \item Symmetric cryptography
    \item Public key cryptosystems
    \item Signature schemes
    \item Hash functions and KDFs
    \item PKI
    \item Protocols \& esoterics
  \end{enumerate}

  Alternatively, there will be limited slots (likely no more than 10) for an oral exam, which covers the same topics but has no T/F section.
  While the oral exam is an alternative to the written one, the written exam is the main expected form of final assessment.

  The written exam lasts 2 hours and the oral exam lasts 30 minutes.
  Both exams are closed book, but a single-sided, handwritten A4 cheat sheet is allowed only for the written exam.

  Students are allowed one retake, provided that their first attempt is not the last available exam attempt.
  There will be at least 3 separate exam dates.
  As per university rules, I will announce the exam dates at least a month in advance, and the last attempt's grade counts.

  The first written exam attempt will be in week 16, with the other two in January.
  The oral exam option will only be available before the winter break.
}

\field{examcriteria}{
  Each T/F question is worth one point.
  Each open-ended question states its point value next to the question, and each of the six sections is worth at least 9 points.
  To pass the exam, the student must score at least 3 points in each of the open-ended sections.
  The course cannot be passed without passing the exam.

  Students who attend at least 8 lectures and 8 practicals earn a 3-point shortfall allowance across the sections.
  For example, the exam is still passed with:
  \begin{itemize}
    \item three sections graded 2;
    \item two sections graded 1 and 2;
    \item one section graded 0.
  \end{itemize}
}

\field{finalgrade}{
  The final grade is conditional on the following requirements:
  \begin{itemize}
    \item The student must earn at least 15/30 points from the homework assignments.
    \item The student must earn at least 10/20 points from the practical tasks.
    \item The student must earn at least 3 points in each of the six open-ended sections in the exam, subject to the shortfall allowance described above.
    \item The student must earn at least one point on the midterm.
  \end{itemize}

  Unless all requirements are satisfied, the final grade is 0, regardless of whether the total points amount to at least 51p.

  \gradescale
}

\stdfield{practices}

\stdfield{ai}

\makescheduletitle

\week{1}{}{Introduction}
\week{2}{}{Randomness, one time pad (OTP), stream ciphers}
\week{3}{}{Block ciphers}
\week{4}{}{Hash functions and key derivation functions (KDF)}
\week{5}{}{Hash structures and message authentication codes (MAC)}
\week{6}{}{Diffie-Hellman (DH), elliptic curve cryptography (ECC), key encapsulation mechanisms (KEM)}
\week{7}{}{ElGamal and RSA}
\week{8}{}{Digital signatures and public key infrastructure (PKI)}
\week{9}{}{Certificate validity, eIDAS, and qualified signatures}
\week{10}{}{TLS, WebPKI and certificate transparency (CT)}
\note{The midterm unlocks on the Monday of week 10 and must be attempted by 23:59 local time that Sunday.}
\week{11}{}{Side channels, hardware security}
\week{12}{}{Commitment schemes and zero-knowledge proofs (ZKP)}
\week{13}{}{Secret sharing, MPC}
\week{14}{}{Internet voting and IVXV}
\week{15}{}{Post-quantum cryptography (PQC)}
\week{16}{}{Revision}
\note{The oral exams and the first written exam will also take place this week.}

\end{document}
